\chapter{Introducción}
\label{ch:introduccion}

La gestión de la seguridad informática en entornos universitarios representa uno de los desafíos más complejos de la ingeniería contemporánea, exigiendo un equilibrio riguroso entre la apertura académica y la salvaguarda de infraestructuras críticas.

\section{Contexto Global de la Ciberseguridad en la Educación Superior}
En la era contemporánea, la transformación digital de las instituciones universitarias ha acelerado la adopción de plataformas web, sistemas de gestión académica y repositorios científicos en la nube. Sin embargo, este proceso de digitalización ha expandido de manera crítica la superficie de ataque institucional. Las universidades gestionan activos de altísimo valor: registros de identidad de miles de estudiantes, calificaciones académicas, patentes e investigaciones confidenciales, y sistemas transaccionales financieros de tesorería y matrícula.

El cibercrimen moderno no discrimina entre corporaciones multinacionales y centros educativos. De acuerdo con informes globales de ciberseguridad, el sector educativo es uno de los más atacados a nivel mundial por grupos de ransomware, actores de ciberespionaje y amenazas internas (\textit{Insider Threats}), debido fundamentalmente a la descentralización de sus redes y a la coexistencia de sistemas legados con aplicaciones modernas.

\section{Importancia de la Evaluación Preventiva de Servicios Web}
Tradicionalmente, las organizaciones han implementado mecanismos defensivos reactivos: despliegue de antivirus tradicionales y firewalls perimetrales que solo responden una vez consumada la intrusión. La ingeniería de ciberseguridad moderna exige la adopción de metodologías \textbf{proactivas} y de diseño seguro (\textit{Secure by Design}). Evaluar la seguridad de los servicios web antes del despliegue en producción o durante auditorías periódicas permite identificar vulnerabilidades estructurales y lógicas que escapan a las herramientas automatizadas.

\section{Integración de Metodologías: STRIDE y Cyber Kill Chain}
Para llevar a cabo una evaluación rigurosa y multidimensional en la Universidad Nacional del Altiplano (UNAP), este proyecto articula dos enfoques metodológicos complementarios:
\begin{enumerate}
    \item \textbf{STRIDE (Microsoft):} Metodología de modelado de amenazas (\textit{Threat Modeling}) orientada al análisis arquitectónico del software y los flujos de datos. Responde a la pregunta: \textit{¿Qué fallas o vectores de riesgo existen en el diseño del sistema?}
    \item \textbf{Cyber Kill Chain (Lockheed Martin):} Marco operativo que describe las siete etapas consecutivas de una intrusión cibernética avanzada. Responde a la pregunta: \textit{¿Cómo avanza un adversario paso a paso para comprometer el objetivo?}
\end{enumerate}

La sinergia entre STRIDE y Cyber Kill Chain permite no solo clasificar las amenazas teóricas, sino comprender su viabilidad técnica y proyectar las contramedidas defensivas más eficaces para interrumpir la cadena de ataque en sus fases tempranas.
